\documentclass[10pt,a4paper]{amsart}
\usepackage[margin=19mm]{geometry}
\usepackage{amsmath,amssymb,mathtools}
\usepackage[scr=rsfs,cal=euler]{mathalfa}
\usepackage{xfrac}
\usepackage[unicode,hidelinks,bookmarks=false]{hyperref}
\newcommand{\ie}{i.e.\ }
\newcommand{\cf}{cf.\ }
\newcommand{\resp}{respectively\ }
\newcommand{\Subsection}{\subsection}
\providecommand{\nobreakdash}{\nobreak\hbox{-}}
\setlength{\emergencystretch}{2em}
\multlinegap0pt
\usepackage{amssymb}
\newtheorem{lem}{Lemma}
\def\N{\mathbb{N}}
\def\ds{\displaystyle}
\title{Perturbation method for non-convex variational problems}
\author{H. Belcheva, N. Zlateva}
\date{Source: arXiv:2606.09053v1 (CC BY 4.0)}
\begin{document}
\maketitle

\noindent\emph{Source and licence.} Perturbation method for non-convex variational problems. Version arXiv:2606.09053v1 (CC BY 4.0), distributed by arXiv under the Creative Commons Attribution 4.0 International licence, http://creativecommons.org/licenses/by/4.0/, as recorded in the arXiv licence metadata for this exact version and shown on the abstract page https://arxiv.org/abs/2606.09053v1. Full licence notice: \texttt{licenses/arxiv-2606.09053-CC-BY-4.0.txt}.\par\smallskip

\noindent\textit{Editorial scope.} Counted unit: the article's lem lem:gen-cantor (source0.tex lines 349--370). The article's complete original proof of it is delivered with it (source0.tex lines 371--451, one \texttt{proof} environment of 81 lines). No mathematics is written, repaired or re-derived: the statement and the proof are the article's own lines, in the article's own notation, and no retained line is substituted or re-flowed.  No further source-local context is needed: every cross-reference of every retained line resolves inside the two delivered spans.  Nothing was omitted from any retained span, so the package carries no omission record.  No retained line contains a citation, so the wrapper carries no bibliography and no bibliography entry.  No retained line contains a figure environment, an image-inclusion command or drawing code, so no image is delivered and the package declares no asset dependency.  Editorial formatting only, in addition: an AMS \texttt{amsart} preamble that restates the article's own declarations and macros used by the retained lines, namely the environments \texttt{lem} on separate counters, together with the standard packages those lines need.  The retained environments print in this document as Lemma 1 in order of appearance, whereas the article prints the same items with the numbering of its own full document.  The complete native source of the article as distributed in the arXiv e-print for this exact version is bundled unchanged as \texttt{sources/202375/source0.tex}.
\begin{lem}[\textbf{Generalized Cantor Lemma}]
    \label{lem:gen-cantor}
    Let $(X,\bar\tau)$ be a topological space and let $d$ be a sequentially $\bar\tau$ lower semicontinuous metric on $X$ such that $(X,d)$ is a complete metric space and, moreover,
    \begin{equation}
        \label{eq:tops-rshp}
        \forall x\in U \in \bar\tau\ \exists \varepsilon > 0,V\in\bar\tau:\ x\in V\text{ and } V_\varepsilon\subset U,
    \end{equation}
    where $V_\varepsilon := \{x\in X:\ d(x,y)\le\varepsilon$ for some $y\in V\}$.

    Let for a set $A\subset X$, $\bar\mu (A)$ denotes the measure of $\bar \tau$ non-compactness of $A$, that is
    $$
        \bar\mu (A) := \inf\{\varepsilon:\ A \subset C_\varepsilon,\text{ for some sequentially }\bar\tau \text{-compact set } C\}.
    $$

    If the sets $A_n\subset X$ are non-empty, nested, i.e. $A_{n+1}\subset A_n$ for all $n\in\mathbb{N}$, sequentially $\bar\tau$-closed, and such that
    $\ds
        \lim_{n\to\infty} \bar\mu(A_n) = 0$,
        then
    $$
        \bigcap_{n=1}^\infty A_n \neq \varnothing.
    $$
\end{lem}

\begin{proof}
    For any $n\in \N$, fix
    $
        x_n \in A_n $.

    Because the sets $A_n$ are nested and sequentially $\bar\tau$  closed, it is enough to show that the sequence $\{x_n\}_{n=1}^\infty$ has a $\bar\tau$-convergent subsequence. To find the latter, we will perform a Cantor diagonalization procedure.

    Let $\varepsilon_n > 0$ be such that
    $$
        \varepsilon_{n+1} < \varepsilon_n,\ \forall n\in\mathbb{N},\text{ and }\lim_{n\to\infty} \varepsilon_n = 0,
    $$
    and for some sequentially $\bar\tau$-compact sets $C_n$, $
        A_n \subset (C_n)_{\varepsilon_n}$.

    Let
    $
        y_n^{(1)} \in C_1$ be such that $d(x_n,y_n^{(1)}) \le \varepsilon_1$.

    Since $C_1$ is sequentially $\bar\tau$-compact, there is $\bar\tau$-convergent subsequence of $\{y_n^{(1)}\}_{n=1}^\infty$. That is, there is strictly increasing function $\sigma_1:\mathbb{N}\to\mathbb{N}$ such that
    $$
         y_{\sigma_1(n)}^{(1)} \to_{\bar\tau} z_1\text{ as } n\to \infty.
    $$
    We choose subsequence of subsequence by induction. Let $\{y_{\sigma_k(n)}^{(k)}\}_{n=1}^\infty\subset C_k$ be such that
    $$
         y_{\sigma_k(n)}^{(k)}\to_{\bar\tau} z_k \text{ as } n\to \infty.
    $$
    Since $\{x_{\sigma_k(n)}\}_{n=1}^\infty$ is a subsequence of $\{x_n\}_{n=1}^\infty$, we have that
    $$
        x_{\sigma_k(n)} \in A_{k+1} \subset (C_{k+1})_{\varepsilon_{k+1}},\quad\forall n \ge k + 1.
    $$
    So, there are
    $$
        \{y_{\sigma_k(n)}^{(k+1)}\}_{n=k+1}^\infty \in C_{k+1}\text{ such that }d(x_{\sigma_k(n)},y_{\sigma_k(n)}^{(k+1)}) \le \varepsilon_{k+1}.
    $$
    Since $C_{k+1}$ is sequentially $\bar\tau$-compact, there is a $\bar\tau$-convergent subsequence of $\{y_{\sigma_k(n)}^{(k+1)}\}_{n=k+1}^\infty$, that is, there is a strictly increasing function $\nu_k$ from $k+1,k+2,\ldots$ into itself such that for
    $
        \sigma_{k+1}(n) := \sigma_k(\nu_k(n))$
    we have that
    $$
         y_{\sigma_{k+1}(n)}^{(k+1)} \to_{\bar\tau} z_{k+1}\text{ as } n\to \infty,
    $$
    and so on \ldots.

    We claim that the sequence $\{z_k\}_{k=1}^\infty$ is  convergent in the metric $d$. Indeed, for any $j>k$, because of the inclusion $\{\sigma_j(n):\ n\in\N\} \subset \{\sigma_k(n):\ n\in\N\}$, we have that
    $$
        d(x_{\sigma_j(n)},y_{\sigma_j(n)}^{(k)}) \le \varepsilon_k.
    $$
    But also, by construction,
    $$
        d(x_{\sigma_j(n)},y_{\sigma_j(n)}^{(j)}) \le \varepsilon_j < \varepsilon_k.
    $$
    So, by the triangle inequality, $d(y_{\sigma_j(n)}^{(k)},y_{\sigma_j(n)}^{(j)}) \le 2 \varepsilon_k$.

    Using that
    $
         y_{\sigma_j(n)}^{(k)} \to_{\bar\tau} z_k$, and $  y_{\sigma_j(n)}^{(j)} \to_{\bar\tau} z_j$, as $n\to \infty$,
        and the $\bar\tau$ lower semicontinuity of the metric $d$, we get
    $$
        d(z_j,z_k) \le 2 \varepsilon_k,\quad\forall j \ge k,
    $$
    which means that the sequence $\{z_k\}_{k=1}^\infty$ is a $d$-Cauchy sequence, and because $(X,d)$ is complete, we get a $z$ such that
    $$
         d(z_k,z) \to 0\text{ as } k\to \infty.
    $$

    Let $U\in\bar\tau$ be an arbitrary neighbourhood of $z$, that is, $z\in U$. Let $V\in\bar\tau$ and $\varepsilon>0$ be such that $z\in V$ and $V_\varepsilon\subset U$, see \eqref{eq:tops-rshp}. Let $j\in\mathbb{N}$ be so large that
    $$
        \varepsilon_j < \varepsilon\text{ and }z_j\in V.
    $$
    Now, since $
        \{\sigma_k(k):k\in\mathbb{N},\ k\ge j\} \subset \{\sigma_j(n): n\in\mathbb{N}\}$,
    and since $y_{\sigma_j(n)}^{(j)} \to_{\bar\tau} z_j$, as $n\to \infty$, there is $N\in\mathbb{N}$ such that
    $$
        y^{(j)}_{\sigma_k(k)} \in V,\quad \forall k > N.
    $$
    But $d(x_{\sigma_k(k)},y^{(j)}_{\sigma_k(k)}) \le \varepsilon_j < \varepsilon$, therefore
    $$
        x_{\sigma_k(k)} \in V_\varepsilon \subset U,\quad \forall k > N.
    $$
    This means that $\{x_{\sigma_k(k)}\}_{k=1}^\infty$ $\bar\tau$-converges to $z$, and the proof is  completed.
\end{proof}

\end{document}
